\chapter{The Problem of Objectivity}
\label{ch:problem}

\begin{chapterguide}
This chapter frames the book's central problem. It distinguishes a mind-independent
world from the human practices by which claims about that world are produced.
It also introduces the four questions that any adequate philosophy of science
must answer: what exists, how evidence supports belief, what scientific
representations mean, and how institutions affect inquiry.
\end{chapterguide}

Science is successful in an ordinary and powerful sense. It predicts eclipses,
builds vaccines, identifies molecules, detects gravitational waves, estimates
the effects of policies, and changes what human beings can do. Yet success alone
does not settle what science knows. A horoscope can sometimes make a vague
prediction that appears correct. A simple model can forecast well while
misrepresenting the mechanism that generates the forecast. An institution can
produce reliable measurements while excluding the people most affected by its
research questions. Scientific success is evidence, but it is not a complete
theory of evidence.

The question of this book is therefore not whether science works. It is how
scientific work can yield objective knowledge when the work is performed by
finite, situated, interested, and fallible beings. We use concepts to identify
objects. We use instruments that embody theories. We compare observations by
statistical models. We inherit categories from institutions and languages. If
these conditions mediate every route to the world, why should the outcome be
anything more than a socially stabilized description?

The answer cannot be that scientists simply look harder. Nor can it be that
scientific claims are objective because a community accepts them. A community
can be mistaken, and a dissenting individual can be right. The answer must
explain how social practices can improve contact with a world that does not
depend on those practices.

\section{Four questions hidden inside \enquote{objectivity}}

The word \emph{objectivity}\index{objectivity} is used for several different
achievements. A measurement may be objective because different observers can
obtain it. A claim may be objective because it is not tailored to one person's
interests. A method may be objective because it makes its assumptions visible.
A theory may be objective because it tracks features of the world. These are
related but not identical.

This book separates four questions.

\begin{enumerate}[label=\textbf{\arabic*.}]
\item \textbf{The ontological question:} what, if anything, exists independently
of human observation, language, and classification? Is reality made only of
observable events, or should science also commit us to atoms, genes, fields,
social structures, or mechanisms?
\item \textbf{The epistemological question:} how can evidence justify a belief
about that reality? What is the role of observation, experiment, probability,
causal intervention, explanation, and criticism?
\item \textbf{The representational question:} what do theories and models
represent? Do they describe entities, mathematical structures, relations,
possibilities, instruments, or only the data they organize?
\item \textbf{The institutional question:} how do communities, incentives,
division of labour, expertise, funding, authority, and exclusion affect the
production of knowledge?
\end{enumerate}

A position can answer one question well and another badly. Constructive
empiricism gives a disciplined account of belief in observables while remaining
agnostic about unobservables. Structural realism identifies an important
continuity across theory change but may leave the nature of structure unclear.
Social epistemology explains why criticism by others matters but cannot by
itself guarantee that a community tracks the world. A complete account must
connect the four questions without reducing them to one another.

\section{The temptation of a single method}

The search for a single scientific method has an understandable motivation.
If there were one valid procedure, objectivity would be easy to define:
apply the procedure and accept its results. But the sciences do not have one
uniform procedure. Astronomers infer from light that cannot be manipulated in a
laboratory. Geologists reconstruct deep time from traces. Epidemiologists
compare populations under ethical constraints. Particle physicists build
detectors whose output depends on sophisticated calibration. Social scientists
study agents who respond to being studied. Mathematical scientists prove
theorems and use simulations to explore spaces too large for analytic
solutions.

This diversity does not imply that science is methodologically anarchic. It
implies that standards must be abstract enough to travel across practices
without becoming so abstract that they say nothing. The shared standard
proposed here is not a recipe but a question:

\begin{principlebox}
\textbf{Constraint question.} Which features of the claim are forced by
independent interactions with the world, and which features are supplied by the
investigator's choices of concepts, models, instruments, and purposes?
\end{principlebox}

The question can be made precise in different ways. In a randomized trial,
independence may come from treatment assignment and outcome measurement. In
climate science, it may come from the convergence of observations, physical
theory, attribution studies, and model ensembles. In a historical science, it
may come from traces that are causally connected to past events and from
independent archives. In each case the objectivity of the claim depends on the
specific constraints that resist convenient reinterpretation.

\section{Mind-independence is not observer-independence}

A reality can be independent of our concepts while being accessible only
through conceptual and technological mediation. A mountain is not created by
the word \enquote{mountain}, but the boundaries of a geological unit depend on
purposes, scale, and theory. The social category of unemployment is not a mere
illusion, but its operational definition is an institutional achievement. An
electron is not invented by a detector, but the detector's output is not
self-interpreting.

This distinction avoids two bad extremes.

\begin{itemize}
\item \textbf{Naive realism} treats perception as a transparent window and
assumes that scientific categories simply copy natural joints.
\item \textbf{Relativism} treats the dependence of descriptions on concepts as
evidence that there is no fact of the matter beyond local schemes.
\end{itemize}

The middle position requires a theory of mediation. Concepts, instruments, and
models are not obstacles that must be removed before inquiry begins. They are
means of contact. Their quality is tested by what they enable investigators to
detect, manipulate, predict, explain, and correct. A microscope does not show
\enquote{reality without mediation}. It creates a controlled relation in which
some features of a specimen become more accessible and other features become
invisible. Objectivity is achieved when the mediation is understood well enough
that the result does not depend on one unexamined viewpoint.

\section{The architecture of the argument}

The argument proceeds in three stages.

First, the major traditions are reconstructed. Logical empiricism correctly
insisted that public evidence and formal clarity matter. Popper correctly
emphasized risk and criticism. Kuhn correctly showed that research is organized
by historically situated paradigms. Lakatos correctly observed that theories are
tested as evolving programmes rather than isolated sentences. Feyerabend
correctly warned against turning a successful history into a rigid rulebook.
Bayesian epistemology correctly clarifies rational belief revision. Realist and
anti-realist traditions each identify real tensions in the relation between
successful theories and invisible entities.

Second, the practices of science are examined: measurement, model building,
explanation, intervention, mechanisms, social criticism, replication, and
uncertainty. These practices reveal that no single relation to truth is
sufficient. A good scientific representation may preserve a structure without
describing every detail. A causal claim requires more than a fitted curve. A
replication can fail because a result was false, because a population changed,
or because the original effect was context-sensitive. Institutional criticism
can increase objectivity without making evidence a product of power alone.

Third, the book develops \emph{Reflexive Constraint Realism} (\RCR), an original
framework with an ontology, an epistemology, inferential rules, and explicit
limits. The framework holds that the world is mind-independent and generative:
it produces regularities, resistances, capacities, and transformations. Science
does not mirror the whole world. It constructs representations that isolate
and track constraints relevant to a question. A claim is objective when its
support is distributed across sufficiently independent routes of inquiry and
when the inquiry exposes the conditions under which the claim would fail.

\begin{figure}[ht]
\centering
\begin{tikzpicture}[
  node distance=4mm and 3mm,
  box/.style={draw=linegray,rounded corners=2pt,fill=pale,align=center,minimum width=1.65cm,minimum height=0.68cm,font=\scriptsize},
  arrow/.style={-{Latex[length=2mm]},draw=teal,thick}
]
\node[box] (world) {Mind-independent\\processes};
\node[box,right=of world] (contact) {Measurement,\\experiment, trace};
\node[box,right=of contact] (model) {Models and\\inferences};
\node[box,right=of model] (claim) {Scoped\\claim};
\node[box,below=of model] (crit) {Criticism,\\replication,\\intervention};
\draw[arrow] (world) -- (contact);
\draw[arrow] (contact) -- (model);
\draw[arrow] (model) -- (claim);
\draw[arrow] (crit) -- (model);
\draw[arrow] (claim) |- (crit);
\end{tikzpicture}
\caption{The basic direction of \RCR: claims are grounded by contact and
tested by further constraint. The arrows are not a single linear method.}
\label{fig:rcr-cycle}
\end{figure}

\section{What would count as an answer?}

An adequate philosophy of science should not merely repeat that science is
fallible. It should specify how fallibility and objectivity coexist. It should
explain why some revisions are improvements rather than merely changes. It
should say when a model's idealization is harmless and when it distorts the
target. It should distinguish a useful forecast from a confirmed causal
mechanism. It should show how values can guide research without determining
empirical conclusions. It should explain why disagreement can be epistemically
productive in one setting and corrosive in another.

The framework developed later makes four commitments.

\begin{enumerate}
\item Reality is not constituted by inquiry, although inquiry can constitute
objects of social practice and categories of measurement.
\item Evidence is relational: its force depends on a claim, an alternative, a
measurement process, and a background model. This does not make evidence
subjective; it makes its conditions explicit.
\item Scientific representations are selective structures. Their truth is
layered rather than all-or-nothing.
\item Objectivity is a property of an inquiry system, not only of an isolated
sentence. A system is objective when it makes error discoverable by agents who
can challenge one another under transparent, world-sensitive constraints.
\end{enumerate}

These commitments are not asserted as axioms at the beginning because their
meaning depends on the history and practice that follow. The next chapters
therefore begin with the traditions that made the problem unavoidable.

\section*{Chapter summary}

The central problem is how a socially and conceptually mediated practice can
produce knowledge of a mind-independent world. Objectivity has ontological,
epistemological, representational, and institutional dimensions. Science has no
single algorithm, but its diverse methods can be evaluated by asking which
parts of a claim are constrained by independent interactions with the world.
The book's proposed answer, Reflexive Constraint Realism, treats objectivity as
stabilized world-tracking under criticism, intervention, and explicit scope
conditions.

\begin{discussion}
\begin{enumerate}
\item Is a measurement objective if all observers agree but the instrument's
calibration is unknown?
\item Can a model be scientifically valuable even if no part of it is literally
true?
\item What would count as an independent route of evidence in a historical
science where controlled intervention is impossible?
\end{enumerate}
\end{discussion}

\begin{furtherreading}
For orientation, compare the historical reconstruction in Kuhn (1962), the
realist debate in van Fraassen (1980) and Chakravartty (2007), the experimental
emphasis of Hacking (1983), and the social account of objectivity in Longino
(1990).
\end{furtherreading}
